% \subsection*{Схемы алгоритмов функций}
% \begin{figure}[H]
% 	\centering
% 	\includegraphics[width=\textwidth, height=\textheight, keepaspectratio]{1.pdf}
% \end{figure}
% \newpage
% \begin{figure}[H]
% 	\centering
% 	\includegraphics[width=\textwidth, height=\textheight, keepaspectratio]{2.pdf}
% \end{figure}
% \newpage
% \begin{figure}[H]
% 	\centering
% 	\includegraphics[width=\textwidth, height=\textheight, keepaspectratio]{3.pdf}
% \end{figure}
% \newpage

% \subsection*{Заголовочный файл для модулей на C++}
% \lstinputlisting{mod.h}

% \subsection*{Модуль на C++}
% \lstinputlisting{mod1.cpp}

% \subsection*{Модуль на Python}
% \lstinputlisting[language=Python]{mod2.py}

% \subsection*{Модуль на C++ (c STL)}
% \lstinputlisting{mod3.cpp}

% \subsection*{Скриншот Графического Пользовательского Интерфейса}
% \begin{figure}[H]
% 	\centering
% 	\includegraphics[width=\textwidth, height=\textheight, keepaspectratio]{screenshot.png}
% \end{figure}
% \noindent
% \indent В правой части (на цветном фоне) отображаются элементы структуры.
% А в левой части сверху вниз: выпадающий список для выбора модуля, поле для ввода/вывода числа и кнопки действий.
% \newpage